% =========================================================
\chapter{C02 --- Composants R, L, C et lois de Kirchhoff}
% =========================================================

\begin{rembox}[title=En bref]
\textbf{De quoi parle ce chapitre ?} Un composant se décrit par une formule simple ; un
\textbf{montage} demande déjà un \textbf{système d'équations}. Les \textbf{lois de Kirchhoff}
(mailles et n\oe uds) permettent de poser toutes ces équations, puis d'en tirer des formules
de simplification : composants équivalents en série/parallèle (R, L, C), théorème de
\textbf{Millman}, théorème de \textbf{Kennelly} (étoile $\leftrightarrow$ triangle). \\
\textbf{Les trois objectifs} : maîtriser les composants de base (R, L, C) avec des analogies
hydrauliques ; décrire un circuit (loi des n\oe uds = conservation des courants, loi des
mailles = conservation de l'énergie) ; simplifier l'analyse (théorèmes clés).
\end{rembox}

\section{Rappel : la résistance}

\slide{c02-03}{Résistance alimentée par une source de courant puis par une source de tension.}

\begin{defnbox}[title=La résistance]
Le composant passif le plus simple. Loi d'Ohm :
\[ u(t) = R\cdot i(t) \qquad\Longleftrightarrow\qquad i(t) = \frac{u(t)}{R} \qquad (R \text{ en } \Omega) \]
Source de courant continue (constante) : $U = R\cdot I$. Source de tension continue :
$I = U/R$.
\end{defnbox}

\begin{goatbox}
L'analogie hydraulique : la tension, c'est la masse d'eau disponible et sa hauteur ; le
courant, c'est la quantité d'eau qui circule ; la résistance, c'est un tuyau étroit (la
section qui freine l'eau). Plus le tuyau est étroit, moins il passe d'eau pour une même
hauteur.
\end{goatbox}

\section{Rappel : le condensateur}

\slide{c02-04}{Le condensateur : $i = C\,du/dt$, analogie du réservoir d'eau.}

\begin{defnbox}[title={Le condensateur (appelé abusivement \og capacité\fg{})}]
\[ q(t) = C\cdot u(t) \qquad i(t) = \frac{dq}{dt} = C\frac{du}{dt} \qquad (C \text{ en F}) \]
Contrairement à la résistance, il \textbf{n'y a pas de loi simple} entre $U$ et $I$ : il faut
dériver ou intégrer.
\end{defnbox}

\begin{exbox}[title={Courant constant $i(t) = I$ dans un condensateur}]
De $\frac{i(t)\,dt}{C} = du$ :
\[ \int\frac{I\,dt}{C} = \int du = \frac{I\,t}{C} + k = u(t)
   \qquad\text{(allure : équation d'une droite)} \]
\[ \int_{t_0}^{t_1}\frac{I\,dt}{C} = \int_{u(t_0)}^{u(t_1)} du
   = \frac{I\,(t_1-t_0)}{C} = u(t_1) - u(t_0) \qquad\text{(parcours)} \]
Un courant constant charge le condensateur \textbf{linéairement} : la tension monte en rampe.
\end{exbox}

\begin{goatbox}
Le condensateur est un réservoir d'eau : la charge $q$ est la masse d'eau
($M(t) = S\cdot h(t)$), la tension est la hauteur $h$, la capacité $C$ est la surface $S$ du
fond. Si on verse un débit constant, la hauteur monte régulièrement ; un réservoir large
(grande $C$) monte lentement.
\end{goatbox}

\section{Rappel : l'inductance}

\slide{c02-05}{L'inductance : $u = L\,di/dt$.}

\begin{defnbox}[title=L'inductance]
\[ u(t) = L\frac{di}{dt} \qquad\Longleftrightarrow\qquad \frac{u(t)\,dt}{L} = di \qquad (L \text{ en H}) \]
\end{defnbox}

\begin{exbox}[title={Tension constante $u(t) = U$ sur une inductance}]
\[ \int\frac{U\,dt}{L} = \int di = \frac{U\,t}{L} + k = i(t) \qquad\text{(une droite)} \]
\[ \int_{t_0}^{t_1}\frac{U\,dt}{L} = \frac{U\,(t_1-t_0)}{L} = i(t_1) - i(t_0) \]
Une tension constante fait monter le courant \textbf{linéairement}.
\end{exbox}

\begin{rembox}[title=Analogie]
L'analogie hydraulique avec l'inductance n'existe pas vraiment (la diapositive montre une
roue à aubes). On se limite à la loi $u = L\,di/dt$. Il existe une analogie plus subtile entre
l'énergie magnétique $\frac12 L\,I^2$ et l'énergie cinétique $\frac12 m\,v^2$.
\end{rembox}

\begin{goatbox}
L'inductance, c'est une lourde roue à eau : elle n'aime pas changer de vitesse. Pour la
faire tourner plus vite (augmenter le courant) il faut pousser (une tension) ; et une fois
lancée, elle continue sur son élan. C'est comme une masse : $L$ joue le rôle de $m$ et le
courant celui de la vitesse.
\end{goatbox}

\section{Les composants passifs : lois fondamentales}

\slide{c02-06}{Le tableau à connaître par c\oe ur.}

\begin{thmbox}[title={Tableau $u = f_1(i)$ et $i = f_2(u)$}]
Pour un dipôle X entre A et B, $u(t) = u_A(t) - u_B(t)$ :
\begin{center}
\renewcommand{\arraystretch}{1.6}
\begin{tabular}{c|c|c}
 & $u(t) = f_1(i(t))$ & $i(t) = f_2(u(t))$ \\ \hline
R & $R\cdot i(t)$ & $\dfrac{u(t)}{R}$ \\
C & $\dfrac1C\displaystyle\int i(t)\,dt + k$ & $C\dfrac{du}{dt}$ \\
L & $L\dfrac{di}{dt}$ & $\dfrac1L\displaystyle\int u(t)\,dt + k$ \\
\end{tabular}
\end{center}
Quand une intégrale est nécessaire sans bornes connues, le résultat est exprimé \textbf{à une
constante près} ($k$).
\end{thmbox}

